% ============================================================================
\section{Introduction}
\label{sec:intro}
% ============================================================================

Relighting outdoor scenes from uncontrolled photo collections is a long-standing challenge in computer vision and graphics.
Unlike studio environments with known illumination, in-the-wild captures exhibit varying sun position, cloud cover, and atmospheric conditions across images.
Recent neural scene representations---particularly Neural Radiance Fields (NeRFs)~\cite{mildenhall2020nerf} and Gaussian Splatting~\cite{kerbl20233dgs}---have enabled photorealistic novel-view synthesis, but most methods bake illumination into the scene representation, preventing post-capture relighting.

Prior works such as LumiGauss~\cite{joaxkal2024lumigauss} addresses this for outdoor scenes by combining 2D Gaussian Splatting (2DGS)~\cite{huang20242dgs} with per-image appearance embeddings and spherical-harmonic (SH) environment maps.
Each image is assigned a learnable SH environment that modulates Gaussian colors, enabling factorization of geometry, albedo, and illumination.
However, the per-image SH representation has important limitations for outdoor relighting:
\begin{itemize}
    \item SH environments are unconstrained and per-image---they do not share structure across images, making it impossible to interpolate to novel, unseen illumination conditions.
    \item Shadows are encoded only implicitly through per-Gaussian SH feature modulation, which cannot produce sharp cast shadows from geometry.
    \item There is no explicit notion of the sun as a directional source, precluding physically-grounded relighting such as time-of-day simulation.
\end{itemize}

We address these limitations with \ours{}, which makes the following contributions:

\begin{enumerate}
    \item \textbf{Explicit directional sun model with scene-global SH lighting} (\S\ref{sec:sun_model}): A physically-motivated lighting decomposition separating direct sunlight (with Lambert shading and an atmospheric colour prior), ambient light, and a shared global sky SH. Sun intensity, colour correction, and ambient are represented as low-order SH functions of the sun direction, reducing lighting parameters to $O(1)$, enforcing physical consistency across views, and enabling evaluation at novel sun positions.

    \item \textbf{Geometry-based shadow computation} (\S\ref{sec:shadows}): Shadow mapping via orthographic sun projection applied to the direct lighting term.

    \item \textbf{Learnable sun direction calibration} (\S\ref{sec:sun_cal}): Per-camera learnable offsets that jointly refine noisy input sun directions during training.

    \item \textbf{Learnable sky-scene distinction} (\S\ref{sec:sky_mask}): A per-Gaussian continuous parameter supervised by sky masks that controls both shadow casting and shadow receiving.
\end{enumerate}